% Source: https://artofproblemsolving.com/community/c3h1341675_angle_measure
% Stable identity: aops-topic-1341675-q01
The base angles are $\angle ABC=\angle ACB=\dfrac{180^\circ-20^\circ}{2}=80^\circ$.

Construct the equilateral triangle $BCE$ with $E$ on the same side of $BC$ as $A$. Since $\angle EBC=60^\circ<80^\circ$, the point $E$ lies inside triangle $ABC$. Both $A$ and $E$ are equidistant from $B$ and $C$, so line $AE$ is the perpendicular bisector of $BC$ and bisects $\angle BAC$; thus $\angle BAE=10^\circ$. Also $\angle ABE=\angle ABC-\angle EBC=80^\circ-60^\circ=20^\circ$.

\begin{asy}
usepackage("amsmath");
size(4.2cm);
defaultpen(fontsize(10pt));
real s=10*Sin(10), h=10*Cos(10);
pair A=(0,h), B=(-s,0), C=(s,0);
real bc=2*s;
pair D=A+bc*unit(C-A);
pair Eq=(0,bc*sqrt(3)/2);
draw(A--B--C--cycle,linewidth(0.9));
draw(B--D,red);
draw(B--Eq--C,blue+dashed); draw(A--Eq,blue+dashed);
dot(A);dot(B);dot(C);dot(D);dot(Eq);
label("$A$",A,N); label("$B$",B,SW); label("$C$",C,SE); label("$D$",D,E);
label("$E$",Eq,E,blue+fontsize(8pt));
draw(arc(A,1.6,270-10,270+10)); draw(A+1.6*dir(262)--A+(-1.6,-1.2),gray+linewidth(0.4)); label("$20^\circ$",A+(-1.6,-1.2),W,fontsize(9pt));
label("\mbox{}",(min(currentpicture,user=true)+max(currentpicture,user=true))/2,black); // keeps the text after the figure black when LaTeX typesets the labels inline
\end{asy}

Compare triangles $ABE$ and $BAD$:
\begin{itemize}
\item $AB=BA$ (common side);
\item $BE=BC=AD$;
\item $\angle ABE=20^\circ=\angle BAD$.
\end{itemize}
By SAS, $\triangle ABE\cong\triangle BAD$, so $\angle ABD=\angle BAE=10^\circ$. Therefore
\[
\angle DBC=\angle ABC-\angle ABD=80^\circ-10^\circ=\boxed{70^\circ}.
\]
